\chapter{The Reference Environment and the Running Example}
\label{app:refenv}

Every version-specific claim in this book is written against one stack, one
cluster shape, and one running example. This appendix fixes all three, plus
the columnar-format background Part III assumes.

\section{The stack}
\label{sec:refenv-stack}

\Cref{tab:refenv-stack} is the reference environment. A claim elsewhere in
the book that depends on version-specific behavior --- an AQE default, a
config key introduced or renamed in a particular release --- is written
against exactly this stack; where behavior differs meaningfully on an older
or newer version, the text says so explicitly rather than leaving it
implicit.

\begin{table}[htbp]
\centering
\caption{The reference stack}
\label{tab:refenv-stack}
\begin{tabular}{@{}ll@{}}
\toprule
\textbf{Component} & \textbf{Version / detail} \\
\midrule
Apache Spark & 3.5.2 \\
Table format & Apache Iceberg, format version 2 \\
Catalog & Hive Metastore (\Cref{sec:cat-hms}) \\
SQL gateway & Apache Kyuubi (\Cref{sec:tuning-kyuubi}) \\
Cluster manager & Google Kubernetes Engine (GKE) \\
Object storage & Google Cloud Storage (GCS) \\
JVM & OpenJDK 17 \\
File format & Apache Parquet \\
\bottomrule
\end{tabular}
\end{table}

\section{The cluster shape}
\label{sec:refenv-cluster}

Worked numbers throughout the book --- executor counts, task-slot arithmetic,
memory sizing --- assume one concrete shape unless a section says otherwise:
executors with 4 cores and 16~GB of executor memory each (plus the memory
overhead \Cref{sec:executor-sizing} derives separately), running as pods on
GKE with no co-located storage, since GCS is remote object storage rather
than node-local disk (\Cref{sec:sched-locality-speculation}'s locality
discussion depends on exactly this separation). A cluster of 10 such
executors gives 40 task slots --- the number \Cref{sec:shuffle-sizing}'s
shuffle-partition-count reasoning starts from.

\section{The running example: \texttt{trips}}
\label{sec:refenv-trips}

The specimen dataset is a city bus system's vehicle position pings. The fact
table, \texttt{trips}, is an Iceberg table partitioned by
\texttt{day(event\_ts)} (\Cref{sec:ice-partitioning}), holding a few hundred
million rows per month, with the columns in \Cref{tab:refenv-trips-schema}.

\begin{table}[htbp]
\centering
\caption{\texttt{trips} schema}
\label{tab:refenv-trips-schema}
\begin{tabular}{@{}lll@{}}
\toprule
\textbf{Column} & \textbf{Type} & \textbf{Meaning} \\
\midrule
\texttt{trip\_id} & \texttt{string} & Unique identifier for one scheduled trip \\
\texttt{event\_ts} & \texttt{timestamp} & When this position ping was recorded \\
\texttt{route\_id} & \texttt{string} & Foreign key to \texttt{dim\_route} \\
\texttt{vehicle\_id} & \texttt{string} & Foreign key to \texttt{dim\_vehicle} \\
\texttt{arrival\_delay\_s} & \texttt{int} & Seconds late (negative if early) at
  the next stop \\
\texttt{distance} & \texttt{double} & Distance covered since the previous
  ping, in km \\
\bottomrule
\end{tabular}
\end{table}

Three dimensions complete the star schema: \texttt{dim\_route} (\texttt{route\_id},
\texttt{route\_name}, and route-level attributes such as the zone it
serves); \texttt{dim\_vehicle} (\texttt{vehicle\_id}, vehicle type and
capacity); and \texttt{dim\_date} (\texttt{date}, \texttt{year},
\texttt{month}, \texttt{day\_of\_week}), used throughout
\Cref{ch:pruning}'s dynamic-partition-pruning examples to filter
\texttt{trips} indirectly by calendar attributes that are not columns on the
fact table itself. \texttt{dim\_route} is small enough to broadcast under
default settings (\Cref{sec:joins-selection}); \texttt{dim\_date} and
\texttt{dim\_vehicle} are smaller still.

\section{A columnar-format primer}
\label{sec:refenv-columnar}

Part III assumes familiarity with two ideas Parquet is built on, summarized
here for a reader coming from row-oriented formats.

A \textbf{row-oriented} format (CSV, JSON, Avro used row-wise) stores a
complete row contiguously --- reading one column out of a wide table still
means reading every other column's bytes for every row along the way. A
\textbf{columnar} format stores one column's values contiguously across many
rows instead, so reading a subset of columns reads only their bytes
(\Cref{sec:pruning-column}'s \texttt{ReadSchema} pruning depends on exactly
this layout to be worth anything) and values of the same type sit next to
each other, which compresses far better than mixed-type row bytes and is
what a vectorized reader (\Cref{sec:ice-read-vectorized}) processes
efficiently.

Parquet organizes a file as a set of \textbf{row groups} (a horizontal slice
of rows, typically sized to match the input-partition target from
\Cref{sec:shape-jobs}), each holding one \textbf{column chunk} per column,
itself split into \textbf{pages} that are individually encoded and
compressed. Every row group's footer carries per-column \textbf{statistics}
--- min, max, null count --- which is exactly what \Cref{sec:c-pushdown}'s
row-group-level pushdown filters against, and what
\Cref{sec:ice-read-metrics-modes}'s Iceberg-level metrics modes summarize one
level up, at the file level, in the manifest.

\begin{note}
"Compressed" and "encoded" are distinct steps applied in sequence: encoding
(dictionary encoding, run-length encoding) exploits repeated or low-cardinality
values before general-purpose compression (Snappy, Zstd, Gzip) runs on the
result. This is why \Cref{sec:mem-heap}'s decompressed-size warning matters:
a compression ratio of 3--5$\times$ on typical analytic data means an
on-disk size is not a reliable proxy for the deserialized, in-memory size
Spark actually has to hold.
\end{note}
